%======================================================================
\chapter{Introduction}\label{ch:intro}
%======================================================================

\section{What this volume is about}

The classification of homogeneous plane continua (\vone{Theorem~0.1}) splits into an indecomposable and a decomposable case. This
volume treats the indecomposable case:

\begin{theorem}[Theorem~\ref{thm:indec-pseudoarc}]
A nondegenerate plane continuum is homogeneous and indecomposable if and only if it is a pseudo-arc.
\end{theorem}

The ``if'' direction is Bing's (Volume~2). The ``only if'' direction has three steps, each a theorem from the literature:
\begin{enumerate}
\item \textbf{Hereditary indecomposability} (Hagopian 1976; Theorem~\ref{thm:HI}): every homogeneous indecomposable plane continuum is hereditarily
indecomposable. The proof used here is Hagopian's of 1984 (following a plan of Jones): if not, a monotone decomposition produces a homogeneous plane
continuum that is indecomposable and contains an arc, and no such continuum exists (Bing 1960, Jones 1975; Theorem~\ref{thm:noarc}).
\item \textbf{Span zero} (Theorem~C of Hoehn--Oversteegen, from Oversteegen--Tymchatyn 1982; Theorem~\ref{thm:theoremC}): every homogeneous indecomposable plane
continuum has span zero.
\item \textbf{Snake-likeness} (Hoehn--Oversteegen 2016, Theorem~1; Theorem~\ref{thm:ho1}): a hereditarily indecomposable plane continuum of span zero is snake-like.
\end{enumerate}
A nondegenerate snake-like hereditarily indecomposable continuum is a pseudo-arc by definition (\vone{Definition~2.14}).

The volume has ten chapters.

\begin{itemize}
\item \textbf{Chapter~\ref{ch:homog}: First consequences of homogeneity in the plane.} Composants, Effros neighbourhoods, Moore's triod theorem; such a continuum is
atriodic and hereditarily unicoherent, and no proper subcontinuum separates the plane (Hagopian 1975, 1976).
\item \textbf{Chapter~\ref{ch:typelambda}: Continua of type $\lambda$.} Irreducible continua with a monotone map onto $\I$ whose layers have empty interior (Kuratowski \S48);
their existence in a decomposable subcontinuum, and Hagopian's Lemma~3 of 1976.
\item \textbf{Chapter~\ref{ch:decomp}: Hagopian's decomposition.} If $M$ is not hereditarily indecomposable, a monotone map of $M$ onto a homogeneous continuum containing an arc
(Hagopian 1984); proofs in outline.
\item \textbf{Chapter~\ref{ch:noarc}: No arc, and hereditary indecomposability.} Jones's arc components and lines, a recurrence argument, a trapping argument in a collar;
Hagopian's theorem.
\item \textbf{Chapter~\ref{ch:span}: Span zero.} Span, Cook's confluence theorem, Davis's theorem, and Oversteegen--Tymchatyn's Theorem~3.6 and Corollary~3.5.
\item \textbf{Chapter~\ref{ch:strips}: The strip theorem.} Oversteegen--Tymchatyn's Theorem~3.2, by a parity argument for closed polygons.
\item \textbf{Chapter~\ref{ch:snake}: Snake-likeness: the plan of Hoehn and Oversteegen.} Theorem~1 from Theorem~20, via Theorem~4 and grid approximations.
\item \textbf{Chapter~\ref{ch:folds}: Simple folds.} The Krasinkiewicz--Minc lemma, Theorem~8 (maps lift through folds), PL graphs, Lemma~6, Proposition~7.
\item \textbf{Chapter~\ref{ch:stairwells}: Stairwells, and Theorem~20 from Theorem~15.} Sides, straight sets, stairwell structures, Propositions~10--19, Lemma~18.
\item \textbf{Chapter~\ref{ch:ladders}: Theorem~15: ladders.} A separator contains a stairwell of odd height, by an explicit construction from a sign function on a grid.
\end{itemize}

The main results are:

\begin{theorem}[Hagopian 1975--1976; Theorems~\ref{thm:hagopian-lemma1} and~\ref{thm:nosep}]
A homogeneous indecomposable plane continuum is atriodic and hereditarily unicoherent, and none of its proper subcontinua separates the plane.
\end{theorem}

\begin{theorem}[Moore's triod theorem; Theorem~\ref{thm:moore-triod}]
Every family of pairwise disjoint triods in the plane is countable.
\end{theorem}

\begin{theorem}[Hagopian 1984; Theorem~\ref{thm:h84}]
A homogeneous indecomposable plane continuum that is not hereditarily indecomposable has a monotone map onto a homogeneous continuum containing an arc, with
proper fibres.
\end{theorem}

\begin{theorem}[Bing 1960, Jones 1975; Theorem~\ref{thm:noarc}]
No homogeneous indecomposable plane continuum contains an arc.
\end{theorem}

\begin{theorem}[Hagopian 1976; Theorem~\ref{thm:HI}]
Every homogeneous indecomposable plane continuum is hereditarily indecomposable.
\end{theorem}

\begin{theorem}[Theorem~C; Theorem~\ref{thm:theoremC}]
Every homogeneous indecomposable plane continuum has span zero.
\end{theorem}

\begin{theorem}[Hoehn--Oversteegen, Theorem~1; Theorem~\ref{thm:ho1}]
A hereditarily indecomposable plane continuum with span zero is snake-like.
\end{theorem}

\begin{theorem}[Hoehn--Oversteegen, Theorem~20; Theorem~\ref{thm:ho20}]
Let $X$ be a hereditarily indecomposable compactum, $f\colon X\to G$ a map to a PL graph, $M\subseteq G\times(0,1)$ a separator of $G\times\set0$ from $G\times\set1$ in $G\times\I$, and $U\supseteq M$
open. Then $f$ is approximated arbitrarily well by $\pi\circ h$ for maps $h\colon X\to U\cap(G\times\I)$.
\end{theorem}

\section{Where these results are used}

Theorem~\ref{thm:indec-pseudoarc} is the indecomposable case of the classification, assembled in Volume~8. Several tools are used again in the decomposable case
(Volume~7): the facts of Chapter~\ref{ch:homog} (atriodicity, hereditary unicoherence, Effros neighbourhoods), type $\lambda$ continua, and the no-arc theorem in
the form of Theorem~\ref{thm:noarc}.

\section{Where the formal proof differs from the literature}\label{sec:intro-diff}

\begin{enumerate}
\item \textbf{Hagopian's papers.} The topological invariance of plane separation, which Hagopian uses without statement, is replaced by \vfour{Theorem~5.22} in the proof that no
proper subcontinuum separates the plane (Theorem~\ref{thm:nosep}). Two theorems quoted from the literature on irreducible continua are false as cited (Remarks~\ref{rem:end-sets-false} and~\ref{rem:thomas21});
the formal development proves counterexamples and uses correct forms. Several of Hagopian's 1984 lemmas are proved by different routes (Chapter~\ref{ch:decomp}).
\item \textbf{Jones's no-arc theorem.} Jones's property (f) cannot be proved from what his argument uses (the dyadic solenoid satisfies all of it), and is replaced by two
properties (R) and (O); a gap in his density argument is replaced by a recurrence argument; the final planar contradiction is a new trapping argument
(Chapter~\ref{ch:noarc}).
\item \textbf{Oversteegen and Tymchatyn.} Several steps of the proofs of Corollary~3.5 and Theorem~3.2 are unargued or impossible as printed (Chapter~\ref{ch:span}). The strip
theorem is proved, in a slightly weaker form that suffices, by a parity argument for closed polygons, without the Jordan curve theorem (Chapter~\ref{ch:strips}).
\item \textbf{Hoehn and Oversteegen.} Theorem~1 as restated on their p.~211, and Theorem~$1'$, are false for a one-point continuum (Remark~\ref{rem:ho1-point}). Folds are built
as subsets of $E\times\R$ rather than as quotients (Chapter~\ref{ch:folds}). Theorem~8 uses a linear push towards fan centres rather than Tietze extensions; the fans must
be disjoint, which is not said. Proposition~7 is proved without boundary bumping and needs a lemma that is not stated. In Proposition~10 a neighbourhood must be shrunk. The
remarks on pp.~193--194 need connectedness. One sentence of the proof of Proposition~19 is false (Remark~\ref{rem:p208}). Claim~15.1 and the irreducible separator in the proof of
Theorem~15 are replaced by ladders (Chapter~\ref{ch:ladders}).
\item \textbf{Proofs in outline.} Chapter~\ref{ch:decomp} and parts of Chapters~\ref{ch:stairwells} and~\ref{ch:ladders} give long case analyses in outline, marked as such. The complete
arguments are in the formal development, at the declarations named in the ``Formal counterpart'' lines.
\end{enumerate}

\section{Prerequisites and conventions}\label{sec:intro-conv}

The reader should know the material of Volume~3 (irreducible continua, the cut-wire theorem, boundary bumping, upper semicontinuous decompositions,
snake-like continua and the $\varepsilon$-map criterion), the Effros theorem (Volume~3), composants (Volume~2), the plane topology of Volume~4 (crossing parity of polygons, Janiszewski's
theorem) and Volume~5 (the Jordan curve theorem, monotone decompositions of the plane, collars).

\begin{itemize}
\item \textbf{Continua.} A continuum is compact, connected and metric; it may be a point unless ``nondegenerate'' is said.
\item \textbf{The plane.} As in Volumes~4 and~5, $\R^2$ carries the sup metric, so that balls are squares.
\item \textbf{Graphs.} In Chapters~\ref{ch:snake}--\ref{ch:ladders}, $E$ is a real normed space, $G\subseteq E$ a PL graph, $G^\star=G\times\I$, and $\pi$ the first projection of $E\times\R$; $\rfr A$ is
the relative frontier of $A$ in the ambient graph.
\item \textbf{Span.} ``Span zero'' means span zero in the sense of Hoehn and Oversteegen (Definition~\ref{def:span}).
\end{itemize}

\paragraph{Formal counterparts.} As in the earlier volumes, many results are followed by a line ``Formal counterpart:'' naming the corresponding declaration of the
Lean development (all names begin with \lean{HPC.}). These lines may be skipped. Appendix~\ref{app:formal} collects them.

\paragraph{Exercises.} Each chapter ends with exercises. None of them is used later.
